\section{Background and Related Concepts}
\label{sec:background}

Here we review
  packages and publishing models,
  associated security risks,
  mitigation by reproducible and verifiable builds,
  and
  the challenges of verifying artifacts in open-source package ecosystems.

\begin{figure}
\centering
\includegraphics[width=0.9\linewidth]{figures/reproducible-builds-overview.drawio.png}
\caption{
\textbf{Reproducible} builds ensure the same source code produces identical artifacts.
\textbf{Verifiable} builds require an independent party to recover the source revision, build environment, and build instructions needed to reproduce that artifact.
We study the artifact verifiability in decentralized-build ecosystems.}
\label{fig:repro-builds-overview}
\end{figure}


% WHat are package registries? How are artifacts packaged, distributed and used?
% Talk about things like pre-built scripts and post-build scripts and impact on the packaging process
% What is the installation process like?

% WHat are software packages, package registries and impact on risks in the software supply chain?
% How are they created and distributed
% What are the threats in this model and how severe are they?


% \comment[Oreofe]{The artifact format differs by ecosystem: PyPI produces \texttt{.whl} (zip) and \texttt{.tar.gz} (sdist) archives; npm produces \texttt{.tgz} tarballs with a \texttt{package/} prefix; crates.io produces \texttt{.crate} archives (gzipped tarballs) with a \texttt{crate-version/} prefix. \toolname strips these ecosystem-specific prefixes before content-level comparison.}

% \PA{Include the kinds of artifacts each ecosystem produces, so readers know what a WHEEL is.}


% \comment[KC]{Put Here: Their use has been advocated for by several industry standards and government initiatives\cite{SLSA_The_Linux_Foundation_2023, Souppaya_Scarfone_Dodson_2022, EO_2021}}
\subsection{Package Ecosystems and Artifacts}
\label{subsec:bg-package-ecosystems}

Software \textit{packages} are reusable software units that provide functionality to downstream applications~\cite{decan_empirical_2019}.
A \emph{package artifact} is the concrete file published for a package version and installed by its users. It contains the package's code, either source directly, or code transformed through compilation, bundling, transpilation, or minification, along with metadata describing the artifact. A \emph{package ecosystem} comprises the packages for a given language or system, the maintainers who publish them, and the tools used to install them~\cite{decan_empirical_2019}.
% It may contain source code directly, code transformed through compilation, bundling, transpilation, or minification, and metadata describing the artifact.
% Together, the packages for a system or language, the maintainers who publish them, and the tools and infrastructure used to build, distribute, and install them form a package \textit{ecosystem}~\cite{decan_empirical_2019}.

Ecosystems differ in where artifacts are built. In \emph{centralized-build} ecosystems, typically operating-system distributions, ecosystem-managed infrastructure builds artifacts from submitted source packages---Debian, for example, builds \texttt{.deb} artifacts from signed source packages on its own daemons~\cite{noauthor_debian_nodate}. In \emph{decentralized-build} ecosystems such as npm~\cite{noauthor_npm_nodate} and PyPI~\cite{noauthor_pypi_nodate}, maintainers build artifacts themselves---on their own machines, CI/CD services, or third-party infrastructure---and upload the finished artifact. Because the build runs outside ecosystem-controlled infrastructure, the registry retains no authoritative record of how each artifact was produced~\cite{benedetti_empirical_2025}.

\PA{I removed the paragraph about source-oriented ecosystems since the title no longer uses it. We can also bring back the definition of package ecosystems since it is a key term.}
% Package ecosystems differ in how closely distributed artifacts correspond to source code.
% In \textit{binary-oriented} ecosystems, such as Debian~\cite{noauthor_debian_nodate} and Maven~\cite{noauthor_maven_nodate}, artifacts primarily contain compiled or binary representations of source code.
% In \textit{source-oriented} ecosystems, such as NPM~\cite{noauthor_npm_nodate} and PyPI~\cite{noauthor_pypi_nodate}, artifacts more directly preserve higher-level source files, although they may also include generated files, minified code, bundled dependencies, or native extensions.

% \JD{Suggest we convert the following paragraph into a Table (resizebox or font cheats for space) that becomes a referenceable definition by us and readers}

% Package ecosystems also differ in how artifacts are built for distribution.
% In \textit{centralized build} ecosystems, often used by operating system distributions, ecosystem-managed infrastructure builds artifacts from submitted source packages.
% For example, Debian maintainers upload signed source packages containing source code, metadata, and build scripts, and Debian build daemons produce the binary \texttt{.deb} artifacts~\cite{managing_debian_packages}.
% In \textit{decentralized build} ecosystems, such as NPM and PyPI, maintainers build artifacts using their own machines, CI/CD services, or other third-party infrastructure before uploading them to the registry.
% As a result, while decentralized ecosystems provide flexibility, the diversity of build tools, release pipelines, and packaging practices make enforcement of any standards more challenging~\cite{benedetti_empirical_2025}.


\subsection{Package Build and Publishing Models and Security Risks}
\label{subsec:bg-build-publish-workflows}


Package ecosystems define the metadata and tooling maintainers use to build and publish artifacts. 
Most require a package manifest in the source directory specifying the name, version, dependencies, and build configuration (\texttt{Cargo.toml} for Crates.io, \texttt{package.json} for npm, \texttt{pyproject.toml} for PyPI), and a \texttt{.gemspec} for RubyGems; in monorepos, the subdirectory holding the manifest identifies the package. 
Ecosystems also provide canonical build and packaging interfaces while permitting third-party build systems and project scripts, such as \texttt{npm pack} for npm or \texttt{python -m build} for PyPI.
To publish, maintainers can either upload pre-built artifacts to registries or publish from CI/CD workflows via trusted publishing~\cite{noauthor_trusted_nodate, crates-trusted-publishing}.


These models expose two security risks~\cite{noauthor_supply_nodate, ohm_backstabbers_2020, duan_towards_2021}. 
In a \emph{build compromise}, an attacker subverts the maintainer's build tools, dependencies, scripts, or CI/CD to inject malicious code during artifact creation; in a \emph{publishing compromise}, an attacker obtains credentials or tokens and uploads a malicious artifact directly. 
Recent npm and PyPI incidents exhibit both~\cite{stepsecurityDurabletask2026, snykTanstack2026, snyk-litellm-2026, noauthor_account_nodate, axiosPostmortem2026, pandyaMaliciousDYdXPackages2026} and in each the distributed artifact diverges from the reviewed source.


% Decentralized-build ecosystems support two publishing channels: maintainers may upload pre-built artifacts directly using registry credentials or long-lived tokens, or publish from CI/CD workflows via trusted publishing, which binds short-lived credentials to the workflow identity~\cite{noauthor_trusted_nodate, noauthor_cratesio_nodate}. These models expose two
% security risks~\cite{noauthor_supply_nodate, ohm_backstabbers_2020, duan_towards_2021}. In a \emph{build compromise}, an attacker subverts the maintainer's build tools, dependencies, scripts, or CI/CD to inject malicious
% code during artifact creation; in a \emph{publishing compromise}, an attacker obtains credentials or tokens and uploads a malicious artifact directly. Recent npm and PyPI incidents exhibit both~\cite{snyk-litellm-2026, ReproducibleBuildsOrg, noauthor_account_nodate,  stepsecurityDurabletask2026} and in each the distributed artifact diverges from the reviewed source.


% Package ecosystems define the metadata and tooling maintainers use to build and publish package artifacts.
% Most require a \textit{package manifest} in the package source directory that specifies metadata such as the package name, version, dependencies, and build configuration.
% These manifest formats are ecosystem-specific, including \texttt{package.json} for NPM, \texttt{pyproject.toml} for PyPI~\cite{python_tool_nodate}, and \texttt{Cargo.toml} for Crates.io.
% In monorepos, where a repository contains multiple packages, the subdirectory containing the manifest typically identifies the package directory.

% Ecosystems also provide standardized or conventional interfaces for producing artifacts, while allowing maintainers to use third-party build systems and project-specific scripts.
% For example, Python packages can be built with the standardized \texttt{python -m build} frontend~\cite{python_tool_nodate}, which invokes the maintainer-specified build backend, while NPM packages commonly use the conventional \texttt{npm run build} that executes project-defined scripts and package artifacts with standardized \texttt{npm pack}~\cite{ScriptsNpmDocs}.

% Ecosystems also define channels for publishing artifacts.
% For instance, centralized ecosystems publish artifacts within ecosystem-managed infrastructure.
% Decentralized ecosystems often support two kinds of channels.
% Maintainers can \textit{directly upload} pre-built artifacts to registries using registry credentials or long-lived tokens, or leverage \textit{trusted publishing} features from CI/CD workflows that use short-lived credentials bound to the CI/CD workflow identity~\cite{noauthor_trusted_nodate, noauthor_cratesio_nodate}.

% These build and publishing models expose different security risks~\cite{noauthor_supply_nodate, ohm_backstabbers_2020, duan_towards_2021}.
% In a build compromise, an attacker compromises the maintainer's build tools, dependencies, scripts, or CI/CD infrastructure, causing malicious code to be inserted during artifact creation.
% In a publishing compromise, an attacker obtains registry credentials or publishing tokens and uploads a malicious artifact directly to the registry.
% Recent attacks, including compromised NPM packages built through malicious third-party GitHub Actions~\cite{snykTanstack2026} and malicious PyPI uploads using compromised maintainer credentials~\cite{snyk-litellm-2026, stepsecurityDurabletask2026, noauthor_account_nodate}, illustrate both risks.
% In these attacks, the artifact distributed to users differ from the reviewed source code.


% Reproducible and verifiable builds provide independent evidence that a software artifact corresponds to its source code.
% A build is \textit{reproducible} if, given the same source code, build environment, and build instructions, an independent party can produce bit-for-bit identical copies of the specified artifacts~\cite{}.
% Verifiability builds on this property: an artifact is verifiable if, given a built artifact, an independent party can recover sufficient information about the original source code and build environment and reproduce an equivalent copy of the original artifact.
% This independent verification can detect artifacts that were compromised during the build or publishing process.
% For this reason, software supply-chain standards~\cite{}, government agencies~\cite{CISANSAODNI2022}, researchers~\cite{lamb_reproducible_2022}, and industry organizations~\cite{BestpracticesSoftwaresupplychainOssscbestpracticesmd} recommend reproducible builds as a mechanism for increasing trust in distributed software artifacts.


\subsection{Reproducible Builds and Artifact Verification}
\label{subsec:bg-repro-builds}
Reproducible and verifiable builds provide independent evidence that a distributed artifact corresponds to its source code. A build is \emph{reproducible} if, given
the same source code, build environment, and build instructions, an independent party can produce a bit-for-bit identical copy of the artifact~\cite{lamb_reproducible_2022, ReproducibleBuildsOrg}.
\emph{Verifiability} adds a recovery requirement: an artifact is verifiable if an independent party, starting only from the published artifact, can recover enough
information about its source and build environment to reproduce an equivalent copy.
The distinction is consequential. Reproducibility \emph{assumes} the source revision, environment, and instructions are known; verifiability requires \emph{recovering} them from the artifact and its metadata. Verifying an artifact therefore rests on two requirements: a build deterministic enough to reproduce, and rebuild information complete enough to drive that reproduction. Because it can detect artifacts altered during build or
publishing, independent verification is recommended as a supply-chain defense by standards bodies~\cite{noauthor_securing_nodate}, government agencies~\cite{CISANSAODNI2022}, researchers~\cite{lamb_reproducible_2022},
and industry~\cite{ReproducibleBuildsOrg}.
% standards bodies=> OpenSSF or CNCF/TAG
% government => CISA, NSA, ODNI, confident
% industy => rbs.org

The determinism requirement is comparatively well understood: prior work has catalogued the common sources of build non-determinism---timestamps, build paths, filesystem ordering, archive metadata, and embedded randomness~\cite{noauthor_variations_nodate, lamb_reproducible_2022, fourne_its_2023, malkaDoesFunctionalPackage2025}---and the reproducible-builds community has built tooling to remove them, such as \texttt{diffoscope} and \path{SOURCE_DATE_EPOCH}, yielding the high reproducibility rates reported in centrally built ecosystems such as Debian~\cite{lamb_reproducible_2022, DebianReproducibleTests} and Nix~\cite{malkaDoesFunctionalPackage2025}. A complementary line of work relaxes byte-for-byte equivalence, accepting artifacts that preserve relevant structure or behavior~\cite{dietrich_levels_2025}.
The recovery requirement, by contrast, is far less settled, and it is where decentralized package ecosystems diverge from centrally built ones.
% \PA{I dont think the next sentences are necessary at this point.}
% Each pattern is partial: attestations are rare (\Cref{sec:rq1}) and may identify only the triggering workflow rather than the exact build, embedded metadata is ecosystem-specific, and heuristics recover
% repositories far more reliably than exact revisions~\cite{PyRadar}.
% % The practical challenges of reproducible builds are well studied.
% Prior work has identified common sources of build non-determinism, including timestamps, build paths, file-system ordering, archive metadata, and randomness~\cite{}.
% In response, the reproducible builds community has developed tools and practices for diagnosing and removing these sources of variation, such as \texttt{diffoscope} for comparing build outputs and \texttt{SOURCE\_DATE\_EPOCH} for normalizing build timestamps.
% These efforts have contributed to high reproducibility rates in controlled ecosystems such as Debian~\cite{} and Nix~\cite{}.
% Prior work has also proposed techniques for identifying benign non-determinism and verifying artifacts at more relaxed equivalence levels, where outputs may differ at the byte level while preserving relevant structure or behavior.


\subsection{The Metadata Recovery Gap in Package Ecosystems}
\label{subsec:bg-verifiability-gap}

Even if build non-determinism is controlled, an independent verifier still needs the \emph{rebuild metadata}---the source revision, build environment, and build instructions---that produced the published artifact.
Centrally managed ecosystems can provide this metadata by construction. 
For example, Debian~\cite{lamb_reproducible_2022} and Nix~\cite{malkaDoesFunctionalPackage2025} derive rebuild metadata from their managed build infrastructure and preserve it as part of the artifact or package specification. 
In decentralized-build ecosystems, however, artifacts are often produced outside the registry by maintainer-controlled CI/CD workflows, leaving registries without a complete record of the build inputs and process.

Several efforts attempt to close this metadata recovery gap. 
Maven Reproducible Central rebuilds Maven Central artifacts using contributor-provided build specification files, while recent work explores automated recovery of build metadata and generation of such specifications~\cite{hassanshahiUnlockingReproducibilityAutomating2025,keshani_aroma_2024}. 
These approaches use heuristics and provenance attestations~\cite{noauthor_provenance_nodate} to recover source commit and CI/CD workflows to recover build commands. 
However, they are tailored to Maven and do not establish whether the recovered metadata is sufficient for artifact-level verification across source-oriented ecosystems.

Google's OSS Rebuild~\cite{google_oss_rebuild} extends rebuild efforts to Crates.io, npm, and PyPI, using similar metadata-recovery heuristics. However, it focuses on supported popular packages and reports using human augmentation when automation fails~\cite{noauthor_introducing_2025}. 
As a result, it remains unclear how far artifacts in decentralized-build ecosystems can be reproduced using only metadata available to an independent verifier.



% and efforts such as Reproducible Central reproduce Maven artifacts from explicit,
% maintainer-provided build specifications.

% The difficulty is structural rather than merely technical: even with build non-determinism fully controlled, a verifier in a decentralized ecosystem has no authoritative account of how an artifact was built. Centrally built ecosystems
% avoid this by construction. Debian and Nix emit rebuild metadata as a byproduct of their managed build infrastructure and retain it alongside the artifact, and efforts such as Reproducible Central reproduce Maven artifacts from explicit,
% maintainer-provided build specifications. In each case the metadata is \emph{recorded at build time} by infrastructure under ecosystem control.

% However, verifiability remains difficult even when build non-determinism can be controlled.
% To verify an artifact, a rebuilder must identify the source revision, build environment, and build instructions that produced the original artifact.
% We refer to this information as \textit{rebuild metadata}.
% Rebuild systems in ecosystems with centralized build infrastructure, such as Debian and NixOS, can often recover this metadata from build records produced during the original build.
% Similarly, efforts such as Reproducible Central rely on explicit build specifications to reproduce Maven artifacts.



% By contrast, open-source package registries such as NPM and PyPI often publish artifacts with fragmented, incomplete, and undocumented rebuild metadata.
% Systems such as Google's OSS Rebuild attempt to infer missing metadata from registry artifacts, package metadata, and source repositories, sometimes augmented by human intervention.
% This inference-based approach limits scalability and coverage, especially outside selected popular packages.
% As a result, supporting verifiable builds across decentralized open-source package registries remains an open problem.
